\section{Implementace simulací}

Tato část popisuje praktickou implementaci čtyř navazujících modelů v prostředí NetLogo. Modely jsou navrženy tak, aby měly podobnou strukturu rozhraní, postupně rostoucí složitost a aby umožňovaly srovnatelné experimenty.

\subsection{Simulace 1: Základní SEIR}

\textbf{Účel:}
\begin{itemize}
\item ověřit základní mechaniku šíření v populaci
\item získat referenční průběh bez dalších zásahů
\end{itemize}

\textbf{Hlavní parametry:}
\begin{itemize}
\item \texttt{num-people} -- velikost populace
\item \texttt{init-infected} -- počet počátečně nakažených
\item \texttt{transmission-prob} -- pravděpodobnost přenosu při kontaktu
\item \texttt{incubation-days} -- délka inkubace
\item \texttt{mean-illness-duration} -- průměrná doba nemoci
\item \texttt{standard-speed} -- rychlost pohybu agentů
\end{itemize}

\textbf{Očekávané chování při změnách parametrů:}
\begin{itemize}
\item vyšší \texttt{transmission-prob} urychluje nástup epidemické vlny
\item vyšší \texttt{num-people} zvyšuje počet kontaktů a absolutní velikost vlny
\item vyšší \texttt{incubation-days} oddaluje nástup infekční fáze
\item vyšší \texttt{standard-speed} zvyšuje promíchání populace a kontaktů
\end{itemize}

\subsection{Simulace 2: SEIR + úmrtnost + kapacita systému}

\textbf{Účel:}
\begin{itemize}
\item rozšířit model o závažné průběhy (úmrtí)
\item zavést první explicitní omezený zdroj (kapacita systému)
\end{itemize}

\textbf{Nové prvky oproti Simulaci 1:}
\begin{itemize}
\item \texttt{fatality-rate} -- základní pravděpodobnost úmrtí
\item kapacitní parametr (např. \texttt{healthcare-capacity}) reprezentující počet současně zvládnutelných vážných případů
\item při překročení kapacity se efektivní úmrtnost zvyšuje
\end{itemize}

\textbf{Očekávané chování při změnách parametrů:}
\begin{itemize}
\item nižší kapacita vede k prudšímu nárůstu úmrtí při vrcholu vlny
\item vyšší \texttt{transmission-prob} zvyšuje riziko překročení kapacity
\item vyšší \texttt{fatality-rate} zvyšuje celkové ztráty i bez přetížení
\end{itemize}

\subsection{Simulace 3: Dočasná imunita, reinfekce a zdroje}

\textbf{Účel:}
\begin{itemize}
\item modelovat dlouhodobou dynamiku (více vln)
\item zavést omezené zdroje využívané jednotlivými agenty
\end{itemize}

\textbf{Nové prvky oproti Simulaci 2:}
\begin{itemize}
\item dočasná imunita po uzdravení (\texttt{immunity-hours})
\item reinfekce po oslabení imunity
\item parametr dostupnosti zdrojů (např. \texttt{resource-stock} nebo \texttt{resource-regeneration})
\item spotřeba zdrojů během nemoci (léčba, izolace, péče)
\end{itemize}

\textbf{Očekávané chování při změnách parametrů:}
\begin{itemize}
\item kratší doba imunity zvyšuje pravděpodobnost další vlny
\item nižší zásoba nebo pomalejší obnova zdrojů prodlužuje období vysokého zatížení
\item při vyčerpání zdrojů roste podíl těžkých průběhů a úmrtí
\end{itemize}

\subsection{Simulace 4: Strategie řízení zdrojů a zásahy}

\textbf{Účel:}
\begin{itemize}
\item porovnat různé strategie alokace omezených zdrojů
\item experimentálně ověřit, které nastavení nejlépe stabilizuje systém
\end{itemize}

\textbf{Nové prvky oproti Simulaci 3:}
\begin{itemize}
\item volba strategie (např. priorita rizikových, priorita prvních případů, rovnoměrná alokace)
\item volitelný zásah (např. očkování části populace)
\item přímé porovnání scénářů při stejném seedu a stejných vstupních parametrech
\end{itemize}

\textbf{Očekávané chování při změnách parametrů:}
\begin{itemize}
\item cílená alokace zdrojů typicky snižuje maxima zatížení systému
\item rovnoměrná alokace bývá robustnější při nejistotě, ale nemusí minimalizovat úmrtí
\item kombinace strategie a očkování zpravidla zkracuje dobu stabilizace
\end{itemize}

\subsection{Doporučená metodika experimentů}

Pro každou simulaci je vhodné zachovat stejný experimentální postup:
\begin{itemize}
\item pro každou kombinaci parametrů spustit více opakování s různými seedy
\item sledovat časové řady stavů S/E/I/R, úmrtí, zatížení systému a stav zdrojů
\item vyhodnocovat průměr, medián, percentily a extrémy
\item porovnávat scénáře podle metrik:
  \begin{itemize}
  \item maximální zatížení systému
  \item celkový počet úmrtí
  \item doba do stabilizace
  \item počet následných vln
  \end{itemize}
\end{itemize}

\subsection{Vazba na cíle práce}

Tato čtveřice modelů odpovídá požadavku replikační studie:
\begin{itemize}
\item začíná jednoduchým ABM modelem
\item postupně přidává realističtější mechanismy
\item umožňuje analyzovat vliv hustoty, přenosu i omezených zdrojů
\item vytváří jasný experimentální rámec pro porovnání variant
\end{itemize}
